\chapter{Bone Marrow Biopsy}

\section*{Overview}
Bone marrow biopsy and aspiration are procedures that obtain bone marrow tissue and cells, usually from the posterior iliac crest, to evaluate blood-forming disorders.

\section*{Alternative Names}
Bone marrow aspirate; trephine biopsy; bone marrow examination

\section*{Principle}
A needle is passed through the cortex into the medullary cavity. Aspiration draws liquid marrow for smears and flow cytometry; a trephine needle captures a solid core for histology, preserving architecture for staging and fibrosis assessment.

\section*{History}
Bone marrow aspiration was introduced in the 1920s and trephine biopsy in the 1950s. They became standard for diagnosing hematologic malignancy and marrow failure.

\section*{Why It Is Done}
Performed for unexplained cytopenias, suspected leukemia, lymphoma, myeloma, or myelodysplasia, to stage hematologic malignancy, to evaluate iron stores and metastatic disease, and to monitor response to therapy.

\section*{Is This a Primary or Secondary Test or Procedure}
Primary diagnostic and staging procedure for many hematologic disorders.

\section*{How It Is Performed}
Performed at the bedside or in clinic with local anesthesia. The patient lies prone or lateral. A Jamshidi-type needle is advanced into the posterior superior iliac spine, aspiration is performed, then a core is obtained with a twisting motion. Specimens are sent for smears, flow cytometry, cytogenetics, and histology.

\section*{Preparation}
Usually none. Informing the clinician about anticoagulants, aspirin, and bleeding disorders is important. Local anesthetic is infiltrated at the site.

\section*{Contraindications and Precautions}
Severe coagulopathy or profound thrombocytopenia increases bleeding risk and may require correction. Osteomyelitis or cellulitis over the site is a contraindication. The sternum is used only for aspiration, never trephine.

\section*{Risks}
Bleeding, infection, and pain are uncommon. Sternal injury is avoided by using the iliac crest. Hematoma may occur at the site.

\section*{Aftercare and Recovery}
Pressure is held over the site for 5-10 minutes, and the patient rests briefly. Mild soreness is treated with acetaminophen or ice.

\section*{Results and Interpretation}
Results reveal cellularity, myeloid-to-erythroid ratio, blast percentage, the presence of malignant cells, fibrosis, and iron stores. Flow cytometry and cytogenetics add immunophenotypic and karyotypic detail.

\section*{Expected Results}
A normocellular marrow (approximately 30-70 percent by age) with normal trilineage maturation and no abnormal infiltrate or fibrosis.

\section*{Follow-up}
Results guide further workup such as cytogenetics, FISH, or molecular panels, and serial biopsies may monitor treatment response.

\section*{References}
\begin{enumerate}
  \item MedlinePlus. Bone marrow biopsy. U.S. National Library of Medicine; 2025.
  \item Current Medical Diagnosis and Treatment. McGraw-Hill; 2025.
\end{enumerate}

\clearpage
